\documentclass{situationist-preprint}

\title{The Ethics of Human Enhancement}
\msid{ETH-22-a}
\msfamily{ETH-22}
\msdiscipline{Normative and Applied Ethics}
\mstype{Principal manuscript}
\mskeywords{human enhancement; therapy; consent; fairness; medical ethics; giftedness; reversal test; positional goods}

\begin{document}
\maketitle

\begin{abstract}
Biomedical interventions can now, or soon will, improve human capacities beyond what is needed for health: memory and attention, muscle strength, resistance to disease, perhaps lifespan. Many hold that such enhancement is wrong, or at least subject to special restrictions, because it goes beyond therapy. We argue that human enhancement is permissible subject to the same constraints of consent and fairness as other medical interventions, no more and no fewer. The argument is a parity argument. The constraints that govern medical interventions attach to features that therapy and enhancement share: risk to the person treated, the possibility of valid consent, and the distribution of access and of effects on others. The feature that distinguishes enhancement, that it raises capacities above normal species functioning, does not attach to any further constraint, since normality has no moral authority of its own. We then examine the objections from giftedness, from authenticity, from the interests of future children and from positional competition, and show that each either reduces to a constraint of consent or fairness or identifies no wrong.
\end{abstract}

\section{The problem}

The debate about enhancement is organized around the distinction between therapy and enhancement. Therapy restores or maintains normal functioning; enhancement improves functioning beyond the normal. On one widely held view, therapy is permissible and its pursuit a central task of medicine, while enhancement is at best tolerated and at worst a threat to human nature, to fairness or to the meaning of human achievement.

The opponents of enhancement have offered several arguments. \citet{Kass2003} and the \citet{PCB2003} report argued that enhancement threatens the dignity of human activity by severing achievements from the efforts that give them their meaning. \citet{Sandel2007} argued that it expresses a drive to mastery that corrodes an appreciation of the giftedness of human powers, and with it the attitudes of humility, solidarity and acceptance on which parenthood and social cooperation depend. \citet{Habermas2003} argued that genetic enhancement of children makes the enhanced person the product of another's design, undermining her capacity to regard herself as the sole author of her life.

Defenders of enhancement have replied that these arguments either prove too much, condemning education and vaccination as well, or rest on a conception of human nature that has no normative authority \citep{Harris2007,Buchanan2011}. \citet{Savulescu2001} went further and argued that parents have reason to select the child with the best expected life. \citet{BostromOrd2006} proposed the reversal test as a diagnostic of status quo bias: if it is held that an increase in some capacity would be bad, one should ask whether a decrease would be good; if neither is, the burden falls on whoever holds that the current level is optimal.

We defend a parity thesis. Enhancement is permissible exactly when an intervention that resembles it in every respect relevant to consent and fairness, but is therapeutic, would be permissible. The thesis is more restrictive than some defenders of enhancement would like, since it subjects enhancement to the full apparatus of medical ethics, and more permissive than its opponents allow, since it subjects it to nothing more.

\section{Framework}

\begin{definition}[Medical intervention]\label{d:med}
A \emph{medical intervention} is a deliberate modification of a person's body or mind by pharmacological, surgical, genetic or other biomedical means.
\end{definition}

\begin{definition}[Therapy and enhancement]\label{d:enh}
A medical intervention is \emph{therapeutic} if it restores or maintains a capacity at or below the level of normal species functioning; it is an \emph{enhancement} if it raises a capacity above that level \citep{Boorse1977,Daniels1985}.
\end{definition}

\begin{principle}[Consent]\label{p:consent}
A medical intervention is permissible only if it is performed with the valid consent of the person treated, informed of its risks and benefits, or, for a person who cannot consent, in accordance with the standards that govern substituted decisions, which protect the person's interests and her future autonomy.
\end{principle}

\begin{principle}[Fairness]\label{p:fair}
A medical intervention is permissible only if the system that provides it distributes access in accordance with the principles of distributive justice, and if its effects on third parties, including effects on their competitive position, are consistent with fair terms of cooperation.
\end{principle}

Principles~\ref{p:consent} and~\ref{p:fair} are the constraints of standard medical ethics, in the form in which they apply to all interventions \citep{BeauchampChildress1979}. Principle~\ref{p:consent} incorporates the requirement of a reasonable balance of benefit and risk, since valid consent cannot be given to an intervention whose risks are disproportionate and that clinicians may not therefore offer. Principle~\ref{p:fair} leaves the content of distributive justice to political philosophy; POL-01 argues that justice is fairness, the principles that would be chosen behind a veil of ignorance.

\begin{principle}[Parity]\label{p:parity}
If two interventions differ only in a feature that is not morally relevant, they are subject to the same moral constraints.
\end{principle}

\section{The argument}

\begin{lemma}\label{l:shared}
The features to which Principles~\ref{p:consent} and~\ref{p:fair} attach, namely risk, benefit, the capacity of the person treated to consent, the distribution of access and the effects on others, are features that enhancements and therapies can share in every combination.
\end{lemma}

\begin{proof}
For each of these features there are therapies and enhancements that have it to the same degree. A drug that raises memory above the normal level in a healthy adult may have the same risk profile and the same mechanism as a drug that restores memory in early dementia. A vaccine, which raises immune function above what an unvaccinated body achieves, and an antibiotic, which restores it, may be given to persons equally able or unable to consent. Access to an enhancement may be distributed justly or unjustly, as may access to therapy. The effects on others of a treatment that gives one competitor an advantage are the same whether the advantage restores a lost capacity or adds a new one.
\end{proof}

\begin{lemma}\label{l:normal}
That an intervention raises a capacity above the level of normal species functioning is not, by itself, a morally relevant feature.
\end{lemma}

\begin{proof}
The level of normal species functioning is a statistical fact about a reference class. Three considerations show that it has no moral authority of its own. First, it varies with the reference class: the normal visual acuity of humans is below that of hawks, and the normal lifespan of humans has doubled in two centuries, without any change in the permissibility of the interventions that changed it. Second, interventions that have always been accepted cross it. Vaccination raises immune function above the normal unvaccinated level, literacy and numeracy raise cognitive capacities above the level of most humans who have lived, and reading glasses, prescribed to restore vision, are equally able to sharpen vision above the norm. Third, the reversal test applies \citep{BostromOrd2006}: those who hold that raising memory above the norm is wrong do not hold that lowering it below the norm, by a corresponding amount, is right; if neither change is good, the norm must be claimed to be optimal, and no argument has been given that it is. A natural endowment, moreover, is not deserved; POL-19 argues that merit is partly a product of luck, and luck is not deserved. The level that the natural lottery fixes has no special claim to be preserved.
\end{proof}

\begin{theorem}\label{t:main}
Human enhancement is permissible subject to the same constraints of consent and fairness as other medical interventions, and to no further constraint.
\end{theorem}

\begin{proof}
Let $E$ be an enhancement. By Lemma~\ref{l:shared}, there is a possible therapeutic intervention $T$ that agrees with $E$ in every feature to which Principles~\ref{p:consent} and~\ref{p:fair} attach. By Lemma~\ref{l:normal}, the feature that distinguishes $E$ from $T$ is not morally relevant. By Principle~\ref{p:parity}, $E$ and $T$ are subject to the same constraints. Since $T$ is subject to Principles~\ref{p:consent} and~\ref{p:fair}, so is $E$; and since no constraint beyond them applies to $T$ in virtue of its being an intervention, none applies to $E$ in virtue of its being an enhancement. So $E$ is permissible if and only if it satisfies the constraints of consent and fairness.
\end{proof}

Theorem~\ref{t:main} does not say that the constraints are satisfied equally easily by enhancements and therapies. An enhancement given to a healthy person may carry the same risk as a therapy given to a sick one, but it offers a different benefit, and the balance that valid consent requires may be struck differently. This is a difference in the inputs to the same constraint, not a different constraint. Similarly, enhancements that confer competitive advantages are more likely than most therapies to engage the second clause of Principle~\ref{p:fair}. Again the principle is the same.

\section{Three cases}

The parity thesis is tested by cases in which the constraints of consent and fairness give determinate verdicts.

\paragraph{Germline editing for disease resistance.} In 2018 it was reported that embryos had been edited to disable the CCR5 gene, with the aim of conferring resistance to HIV, and that twins had been born from them. Resistance to an infection the children did not have is an enhancement in the sense of Definition~\ref{d:enh}. The case was condemned almost universally, and the parity thesis explains why without appeal to the fact that it was an enhancement. The children could not consent, and the standards of substituted decision-making forbid exposing a person to a large and uncharacterized risk, including off-target edits and unknown effects of the edited gene, for a benefit available by safer means, since the transmission of HIV from parent to child can be prevented by other methods. A therapeutic intervention with the same profile of risk and benefit would have been equally impermissible. Conversely, a germline edit with well-characterized effects that protected against a common and serious disease, unavailable by other means, would satisfy the constraints whether it counted as therapy or as enhancement.

\paragraph{Cognitive enhancers for surgeons.} Suppose a drug reliably improves the sustained attention of fatigued surgeons, with low risk. A surgeon who takes it with informed consent satisfies Principle~\ref{p:consent}. Its effects on others are mainly beneficial, since patients are safer, and the competitive advantage it confers on surgeons who take it over those who do not is not one whose distribution matters to fair terms of cooperation in the way that a competitive advantage in sport does. The parity thesis permits it. What it does not permit is an employer's requirement that surgeons take it, since consent obtained under threat of dismissal is not valid in the sense of Principle~\ref{p:consent}, unless the requirement is a reasonable condition of the work in the way that vaccination requirements for clinical staff can be.

\paragraph{Erythropoietin in endurance sport.} Erythropoietin raises the oxygen-carrying capacity of the blood. Taken by an athlete with informed consent, it may satisfy Principle~\ref{p:consent}, though at high doses its risks are serious. It violates Principle~\ref{p:fair} where the rules of the sport forbid it, since the athlete gains an advantage over competitors who keep the rules, and its benefits are almost wholly positional. The parity thesis therefore agrees with the prohibition, and locates its ground in fairness rather than in the fact that the drug enhances. A patient with renal anemia who takes the same drug to restore normal function is not subject to the prohibition, because no competitive effect arises, not because the use is therapeutic; a competing athlete with renal anemia would raise exactly the fairness question that the rules of the sport must settle.

\section{Objections and replies}

\begin{objection}[Giftedness]\label{o:sandel}
Enhancement expresses and reinforces a drive to mastery, and erodes the appreciation of our powers as gifts. This attitude sustains humility, an openness to the unbidden in parenthood, and the solidarity of those who know their fortunes might have been otherwise \citep{Sandel2007}.
\end{objection}

\begin{reply}
The objection concerns an attitude, not an act, and it applies to therapy as much as to enhancement. A parent who treats a child's congenital disease also refuses to accept what was given, and no one thinks the refusal reflects a drive to mastery. If the attitude is objectionable, it is objectionable wherever it appears and should be criticized as a vice, not used to restrict a class of interventions. As for solidarity, the argument depends on the claim that people who can enhance themselves will come to regard their advantages as earned. That is an empirical prediction about attitudes, and the correct response to it, if it is true, is to sustain institutions of solidarity, as Principle~\ref{p:fair} already requires.
\end{reply}

\begin{objection}[Authenticity and achievement]\label{o:kass}
An achievement obtained with the help of an enhancement is less fully the agent's own, and so less valuable \citep{Kass2003}. Achievement is a component of the good (ETH-01), and enhancement threatens it.
\end{objection}

\begin{reply}
The value of an achievement depends on the difficulty of what is achieved relative to the agent's capacities and on the agent's exercise of them. An enhanced pianist who plays a difficult piece achieves less, in this sense, than an unenhanced pianist who plays it with greater effort; but she can achieve more by attempting what is difficult for her. Enhancement shifts the scale on which difficulty is measured; it does not abolish achievement, any more than the invention of the piano abolished the achievement of harpsichordists. Where the comparison between competitors is the point of an activity, as in sport, rules on enhancement are rules of fair competition and fall under Principle~\ref{p:fair}.
\end{reply}

\begin{objection}[The child's authorship]\label{o:habermas}
Genetic enhancement of a child, decided by her parents before birth, makes her the product of another's intention in a way that undermines her standing as the author of her own life \citep{Habermas2003}.
\end{objection}

\begin{reply}
The child cannot consent, and Principle~\ref{p:consent} therefore requires that the decision be made according to the standards of substituted decision-making, which protect her interests and her future autonomy. These standards forbid interventions that narrow the child's options, such as an enhancement fitting her for a single career, and permit those that widen them, such as improved resistance to disease or improved general cognitive capacities. Habermas's concern is the protection of the child's future autonomy, and the constraint of consent, in its form for those who cannot consent, already protects it. The objection identifies a reason to restrict some enhancements of children; it does not identify a constraint beyond consent.
\end{reply}

\begin{objection}[Positional arms race]\label{o:position}
Many of the benefits of enhancement are positional: being taller or more attentive matters because it puts one ahead of others. Widespread enhancement would leave everyone where they started, at greater cost and risk, and the first to enhance would gain at others' expense.
\end{objection}

\begin{reply}
The objection identifies a collective action problem that is real for some enhancements, and it falls squarely under the second clause of Principle~\ref{p:fair}, which concerns the effects of interventions on the competitive position of others. Where an enhancement's benefits are mainly positional, fairness may require regulation, as it does for performance-enhancing drugs in sport. But many enhancements are not positional. Resistance to disease, improved memory in old age and a longer healthy life benefit their bearers whether or not others have them, and the objection gives no reason to restrict these. Unequal access raises the same issue of distributive justice as unequal access to any medical good, and it is answered in the same way \citep{Buchanan2000}. ETH-22-b applies this reply to cognitive enhancement.
\end{reply}

\section{What the result does not establish}

Theorem~\ref{t:main} is a claim about permissibility. It does not establish that any agent ought to enhance herself or her children, and so does not endorse the principle of procreative beneficence \citep{Savulescu2001}. It does not establish that any particular enhancement now available satisfies the constraints; for most pharmacological cognitive enhancers, the evidence of benefit in healthy adults is modest and the evidence about long-term risk incomplete, so that valid consent requires caution \citep{Greely2008}. It does not establish that health systems are obliged to fund enhancements. Fairness constrains the distribution of what is provided, but whether a given enhancement should be publicly provided is a question about the priorities of distributive justice, on which the parity thesis is silent.

\section{Conclusion}

Human enhancement is permissible, subject to the same constraints of consent and fairness as other medical interventions. The features to which those constraints attach are shared by therapy and enhancement, and the feature that separates them, the crossing of a statistical norm, has no moral significance of its own. The objections from giftedness, authenticity, the child's authorship and positional competition either reduce to the constraints of consent and fairness or identify attitudes and outcomes that those constraints, applied with care, already address.

\begin{thebibliography}{99}\small
\bibitem[Beauchamp and Childress(1979)]{BeauchampChildress1979} Beauchamp, T.~L. and Childress, J.~F. (1979). \emph{Principles of Biomedical Ethics}. Oxford University Press.
\bibitem[Boorse(1977)]{Boorse1977} Boorse, C. (1977). Health as a theoretical concept. \emph{Philosophy of Science}, 44(4), 542--573.
\bibitem[Bostrom and Ord(2006)]{BostromOrd2006} Bostrom, N. and Ord, T. (2006). The reversal test: eliminating status quo bias in applied ethics. \emph{Ethics}, 116(4), 656--679.
\bibitem[Buchanan(2011)]{Buchanan2011} Buchanan, A. (2011). \emph{Beyond Humanity? The Ethics of Biomedical Enhancement}. Oxford University Press.
\bibitem[Buchanan et~al.(2000)]{Buchanan2000} Buchanan, A., Brock, D.~W., Daniels, N. and Wikler, D. (2000). \emph{From Chance to Choice: Genetics and Justice}. Cambridge University Press.
\bibitem[Daniels(1985)]{Daniels1985} Daniels, N. (1985). \emph{Just Health Care}. Cambridge University Press.
\bibitem[Greely et~al.(2008)]{Greely2008} Greely, H., Sahakian, B., Harris, J., Kessler, R.~C., Gazzaniga, M., Campbell, P. and Farah, M.~J. (2008). Towards responsible use of cognitive-enhancing drugs by the healthy. \emph{Nature}, 456(7223), 702--705.
\bibitem[Habermas(2003)]{Habermas2003} Habermas, J. (2003). \emph{The Future of Human Nature}. Polity Press.
\bibitem[Harris(2007)]{Harris2007} Harris, J. (2007). \emph{Enhancing Evolution: The Ethical Case for Making Better People}. Princeton University Press.
\bibitem[Kass(2003)]{Kass2003} Kass, L.~R. (2003). Ageless bodies, happy souls: biotechnology and the pursuit of perfection. \emph{The New Atlantis}, 1, 9--28.
\bibitem[President's Council on Bioethics(2003)]{PCB2003} President's Council on Bioethics (2003). \emph{Beyond Therapy: Biotechnology and the Pursuit of Happiness}. Washington, DC.
\bibitem[Sandel(2007)]{Sandel2007} Sandel, M.~J. (2007). \emph{The Case against Perfection: Ethics in the Age of Genetic Engineering}. Harvard University Press.
\bibitem[Savulescu(2001)]{Savulescu2001} Savulescu, J. (2001). Procreative beneficence: why we should select the best children. \emph{Bioethics}, 15(5--6), 413--426.
\end{thebibliography}

\end{document}
